\chapter{External literature results}
\label{chap:external-results}

This chapter records published inputs and the specializations to be checked or
proved internally.  A completed node must retain its source and locator; rows
marked partial or missing are statement-writing obligations, not yet verified
external inputs.  The text distinguishes an explicit external hypothesis from
an internal proof.

The stage-delivery record \texttt{Challenge2.LiteratureInterface} separates
Adams/May, Moss, and BR21 conclusions from computation delivery. Its Moss and
tmf claims refer to the same shared contexts; cobar comparison and tmf model
bindings are not reclassified as literature theorems. This organization does
not supply proofs or replace the source-bearing inputs described below.

\section{Master inventory of external statements}
\label{sec:external-result-master-inventory}

This is a working statement-writing inventory for external mathematics quoted
or consumed by the project, not a stage-ownership or computation ledger.  One
row is one \emph{statement package}:
clauses with one source locator and one consumer are kept together, while
mathematically independent results from the same paper receive separate rows.
The inventory has \textbf{53 packages}: 12 common or foundational packages,
28 inputs for the selected proof route, 10 inputs used only by the two-line or
geometric conclusions, and 3 precise background results recorded for citation
completeness.  A cross-check source does not create a second package.

The count is a bibliographic writing inventory, not a count of fields in any
Lean structure. In particular, foundational results may discharge the fixed
implementation or Challenge~2 foundation obligations, route results may occur
in its literature interface, and citation-only
background need not become a stage input at all.

The status column has the following meaning.  ``Node'' means that this chapter
already has an individual Blueprint node; ``grouped'' means that the statement
is currently visible only as a field of one of the four route input records;
``partial'' means that the existing node is strictly weaker than the quoted
result; and ``missing'' means that an individual statement node still has to be
written.  None of these status words asserts that a proof or a fixed-model
witness has been constructed.

\subsection{Common and foundational results: 12 packages}

\begin{center}
\scriptsize
\begin{tabular}{p{0.07\linewidth}p{0.43\linewidth}p{0.25\linewidth}p{0.14\linewidth}}
\textbf{ID} & \textbf{Statement package} & \textbf{Primary source} & \textbf{Status} \\
\hline
F01 & Classical mod--2 Adams convergence and detection for the stated complete finite-type spectra & \texttt{Adams58} & Node: \ref{thm:external-classical-adams-convergence} \\
F02 & The sphere Adams one-line is generated by the classes $h_j$ in degrees $(1,2^j)$ and vanishes elsewhere & \texttt{Adams58} & Node: \ref{thm:external-adams-one-line} \\
F03 & Hopf-invariant-one survival for $j\le3$ and $d_2(h_j)=h_0h_{j-1}^2\ne0$ for $j\ge4$ & \texttt{Adams60} & Node: \ref{thm:external-adams-one-line} \\
F04 & Adams filtration at least $k$ iff a map factors into $k$ mod--2-homology-zero maps & \texttt{RavenelGreenbook}, Theorem 2.2.14 & Partial: \ref{thm:external-map-filtration-factorization} has only the forward direction \\
F05 & Existence of the stable presentable symmetric-monoidal synthetic category and the synthetic analogue functor $\nu$ & \texttt{Pst} & Node: \ref{thm:external-synthetic-foundation} \\
F06 & $\nu$ carries a cofiber sequence to a cofiber sequence iff mod--2 homology is short exact & \texttt{Pst}, Lemma 4.23 & Node: \ref{thm:external-nu-cofiber-criterion} \\
F07 & The first quotient $S/\lambda$ has the unique compatible commutative algebra structure & \texttt{Pst}, Corollary 4.45; \texttt{BHSmot} & Node: \ref{thm:external-lambda-quotient-ring} \\
F08 & The finite quotients $S/\lambda^n$ form a coherently multiplicative commutative-algebra tower & \texttt{BHSmot}, Appendices B--C; \texttt{BurklundXu}, Construction 7.7 & Grouped in the route algebra input \\
F09 & $\lambda$-inversion identifies synthetic spectra with the classical generic fiber, compatibly with $\nu$ and tensor products & \texttt{Pst} & Node: \ref{thm:external-lambda-inversion} \\
F10 & The $\lambda$-special fiber is the derived category of dual-Steenrod-algebra comodules & \texttt{Pst}; \texttt{Hovey}; checked in \texttt{BHS} & Missing \\
F11 & For an $H\mathbb F_2$-nilpotent-complete spectrum, $\nu X\simeq\varprojlim_n\nu X/\lambda^n$ & \texttt{BHS}, Proposition A.13 & Node: \ref{thm:external-synthetic-lambda-complete} \\
F12 & Pontryagin--Thom identifies stable sphere homotopy with framed cobordism in the convention used for Kervaire classes & Classical Pontryagin--Thom theorem; exact project locator pending & Missing \\
\end{tabular}
\end{center}

\subsection{Selected proof-route inputs: 28 packages}

\begin{center}
\scriptsize
\begin{tabular}{p{0.07\linewidth}p{0.43\linewidth}p{0.25\linewidth}p{0.14\linewidth}}
\textbf{ID} & \textbf{Statement package} & \textbf{Primary source} & \textbf{Status} \\
\hline
A01 & An order-two classical $\theta_5$ detected by $h_5^2$ exists & \texttt{Xu}, Corollary 1.3; \texttt{BJMtheta5}; \texttt{IWX} & Grouped in \ref{def:route-classical-bx-inputs} \\
A02 & Every element of the classical $62$-stem has exponent two & \texttt{IWX}, classical $62$-stem table & Grouped in \ref{def:route-classical-bx-inputs} \\
A03 & Two choices detected by $h_5^2$ differ by a class of Adams filtration at least six & \texttt{IWX}, filtrations $2,6,8,10$ in the $62$-stem & Grouped in \ref{def:route-classical-bx-inputs} \\
A04 & The actual maps $2$, $\eta$, and $\nu$ are detected by the standard $h_0,h_1,h_2$ classes & Adams low stems; \texttt{BHS}; \texttt{IWX} & Grouped in \ref{def:route-classical-bx-inputs} \\
A05 & The original finite and untruncated BJM/BX criterion for one distinguished $\theta_5$ & \texttt{BurklundXu}, Proposition 7.19 & Node: \ref{thm:external-bjm-bx-criterion} \\
A06 & The quadratic-construction bottom/top-cell calculation used for that distinguished choice & \texttt{BurklundXu}, Theorem 7.15 and proof of Proposition 7.19 & Node: \ref{thm:external-bx-quadratic-cell-calculus} \\
A07 & The source-choice identity $\delta_1(h_6^2)=\lambda\eta(\theta_5^{\mathrm{src}})^2$ & \texttt{BurklundXu}, proof of Proposition 7.19 & Node: \ref{thm:external-total-differential-identity} \\
A08 & A positive-filtration classical map admits a full synthetic $\lambda$-factor lift & \texttt{BHS}, Lemma 9.15 & Node: \ref{thm:external-synthetic-lift} \\
A09 & A suitable classical triangle admits the specified triangle-compatible synthetic lift & proof of \texttt{BHS}, Lemma 9.15 & Node: \ref{thm:external-synthetic-triangle-lift} \\
A10 & A class lies in $Z_q$ iff it lifts to the actual finite quotient $\nu X/\lambda^q$ & \texttt{BHS}, Theorem A.1(1a--b) & Node: \ref{thm:external-finite-lambda-bockstein} \\
A11 & A suitable finite lift has Bockstein boundary representing the corresponding Adams differential & \texttt{BHS}, Theorem A.1(1c) & Node: \ref{thm:external-finite-lambda-bockstein} \\
A12 & A common $Z_\infty$ representative lifts to untruncated $\nu X$, and conversely & \texttt{BHS}, Theorem A.1(2) & Grouped in \ref{def:route-synthetic-inputs} \\
A13 & Classical and synthetic Adams differentials correspond, including the required $E_2$ weight-vanishing range & \texttt{BHS}, Theorem A.8 & Node: \ref{thm:external-synthetic-rigidity} \\
A14 & The $E_\infty$ page of $\nu X$ is the stated cycle/boundary quotient & \texttt{BHS}, Corollary A.9 & Node: \ref{thm:external-synthetic-einfty-nu} \\
A15 & The $E_\infty$ page of $\nu X/\lambda^q$ is the stated truncated cycle/boundary quotient & \texttt{BHS}, Corollary A.11 & Node: \ref{thm:external-synthetic-einfty-quotient} \\
A16 & Realization kills a selected synthetic class iff some finite power of $\lambda$ kills it & \texttt{Pst}, localization and completion results & Grouped in \ref{def:route-kernel-source-binding} \\
A17 & Adams filtration is equivalent to the specified $\lambda$-divisibility bound & \texttt{BHS}, \texttt{cor:tau-surj} & Grouped in \ref{def:route-synthetic-inputs} \\
A18 & A nonzero class surviving through $E_{r+1}$ has nonzero $\lambda^{r-1}$ multiple & \texttt{BHS}, Theorem A.1(2a) & Grouped in \ref{def:route-synthetic-inputs} \\
A19 & Permanent classes and prescribed classical detections admit the stated realization-compatible lifts & \texttt{BHS}, Theorem A.1(2b),(3b) & Grouped in \ref{def:route-synthetic-inputs} \\
A20 & If a permanent class is hit by $d_r$, some lift is killed by $\lambda^{r-1}$ & \texttt{BHS}, Theorem A.1(3a) & Grouped in \ref{def:route-synthetic-inputs} \\
A21 & May TC3 supplies the signed smash-boundary/push--pull conclusion with its compatible-pair hypothesis & \texttt{May01}, TC3 and Lemma 4.6 & Node: \ref{thm:external-may-smash-boundary} \\
A22 & Moss convergence supplies a permanent Massey member detecting a corresponding Toda member under all crossing hypotheses & \texttt{Moss}, Theorem 1.2; Belmont--Kong and \texttt{IWX} cross-checks & Grouped in \ref{def:route-moss-source} \\
A23 & The low synthetic Toda relations and the symmetric bracket consequence used in stem $62$ & \texttt{BHS}, \texttt{prop:syn-toda-range}; Toda; \texttt{IWX}, \texttt{cor:2-symmetric} & Grouped in \ref{def:route-toda-source-application} \\
A24 & The classical $\theta_5$ maps to zero under the actual tmf Hurewicz map & \texttt{tmf} (Behrens--Mahowald--Quigley), Theorem 1.2 and Figure 1.1 & Grouped in \ref{def:route-tools-tmf-inputs} \\
A25 & The class $g^4\Delta h_1g$ survives and detects a class with nonzero tmf image in stem $125$ & \texttt{tmf}, Section 7 & Grouped in \ref{def:route-tools-tmf-inputs} \\
A26 & The tmf detector has the required stem-$63$ low-filtration vanishing range & \texttt{tmf}, Section 2 and change of rings & Grouped in \ref{def:route-tools-tmf-inputs} \\
A27 & The intrinsic $\lambda$-Bockstein page and differential agree with the regraded synthetic Adams page and differential & \texttt{BHS}, Theorem A.1 and its Bockstein-page comparison & Node: \ref{thm:external-lambda-bockstein} \\
A28 & In the ordinary Adams spectral sequence of tmf, $d_3(v_2^{16})=\beta^5g$ with the stated coordinates & \texttt{BR21} & Node: \ref{thm:external-tmf-br21-d3-slice} \\
\end{tabular}
\end{center}

\subsection{Two-line and geometric conclusions: 10 packages}

\begin{center}
\scriptsize
\begin{tabular}{p{0.07\linewidth}p{0.43\linewidth}p{0.25\linewidth}p{0.14\linewidth}}
\textbf{ID} & \textbf{Statement package} & \textbf{Primary source} & \textbf{Status} \\
\hline
C01 & $h_0h_2,h_0h_3,h_2h_4$ and $h_j^2$ for $j\le3$ survive nontrivially & \texttt{Maythesis} & Node: \ref{thm:external-adams-one-line} \\
C02 & The Adams four-line calculation and Hopf differentials give the displayed complete $E_3$ candidate list on the two-line & \texttt{LinExt4}; \texttt{Adams60} & Missing \\
C03 & $d_3(h_2h_5)\ne0$, $d_4(h_3h_5)\ne0$, and $h_4^2$ survives & \texttt{MahowaldTangora} & Missing \\
C04 & $h_3h_6$ supports the cited nonzero $d_3$ differential & \texttt{IWX} & Missing \\
C05 & The family $h_1h_j$ survives for every $j\ge3$ & \texttt{Mahowaldetaj} & Missing \\
C06 & Every $h_j^2$ for $j\ge7$ supports a nonzero differential & \texttt{HHR} & Missing; the geometric HHR node below is not this spectral-sequence statement \\
C07 & If $2\theta_j=0$ and $\theta_j^2=0$, an order-two $\theta_{j+1}$ exists & \texttt{BJMinduction} & Node: \ref{thm:external-bjm-induction} \\
C08 & Browder's equivalence between Kervaire-one framed manifolds and nonzero permanence of the corresponding $h_j^2$ & \texttt{Browder} & Node: \ref{thm:external-browder-criterion} \\
C09 & Kervaire-one framed manifolds exist in dimensions $2,6,14,30,62$ & \texttt{Browder}; \texttt{Maythesis}; \texttt{MahowaldTangora}; \texttt{BJMtheta5}; \texttt{Jones30} & Node: \ref{thm:external-low-kervaire-existence}; exact per-dimension locators pending \\
C10 & No Kervaire-one framed manifold exists in dimensions $2^{j+1}-2$ for $j\ge7$ & \texttt{HHR} & Node: \ref{thm:external-hhr-nonexistence} \\
\end{tabular}
\end{center}

The $h_5^2$ permanence/order-two input used by the two-line corollary is A01,
and the new $h_6^2$ permanence theorem is internal to the target paper, so
neither is counted a second time in this subsection.

\subsection{Precise background results: 3 packages}

\begin{center}
\scriptsize
\begin{tabular}{p{0.07\linewidth}p{0.43\linewidth}p{0.25\linewidth}p{0.14\linewidth}}
\textbf{ID} & \textbf{Statement package} & \textbf{Primary source} & \textbf{Status} \\
\hline
B01 & The dimension-$10$ Kervaire PL manifold admits no smooth structure & \texttt{Kervaire} & Missing; citation-audit only \\
B02 & The Kervaire--Milnor classification connects stable framed data with exotic smooth structures on spheres & \texttt{KervaireMilnor} & Missing; citation-audit only \\
B03 & Jones constructs an explicit framed Kervaire-one manifold in dimension $30$ & \texttt{Jones30} & Missing as an individual node; C09 records only existence \\
\end{tabular}
\end{center}

The 53 packages involve 28 source entries when cross-check works and the
not-yet-fixed Pontryagin--Thom locator are counted: \texttt{Adams58},
\texttt{Adams60}, \texttt{RavenelGreenbook}, \texttt{Pst}, \texttt{BHSmot},
\texttt{BurklundXu}, \texttt{Hovey}, \texttt{BHS}, Pontryagin--Thom,
\texttt{Xu}, \texttt{BJMtheta5}, \texttt{IWX}, \texttt{May01},
\texttt{Moss}, Belmont--Kong, Toda, \texttt{tmf}, \texttt{BR21},
\texttt{Maythesis}, \texttt{LinExt4}, \texttt{MahowaldTangora},
\texttt{Mahowaldetaj}, \texttt{HHR}, \texttt{BJMinduction},
\texttt{Browder}, \texttt{Jones30}, \texttt{Kervaire}, and
\texttt{KervaireMilnor}.

Pure definitions, broad historical surveys, and computational artifacts are
not statement packages in this count.  In particular, \texttt{LWXMachine},
\texttt{LWXZenodo}, \texttt{LinPlot}, CSV/proofs.db rows and Appendix tables
belong to computed evidence $C(M)$.  Citations such as \texttt{WX61},
\texttt{IWXsurvey}, \texttt{WXems}, \texttt{WXicm}, \texttt{Brunerdiff},
\texttt{DuggerIsaksen}, \texttt{Isaksen}, \texttt{IsaksenXu},
\texttt{LinalgKP}, \texttt{GWX}, \texttt{Wang},
\texttt{BurklundIsaksenXu}, \texttt{LinGrobner}, \texttt{BJ},
\texttt{Nassau}, \texttt{Behrens}, \texttt{Ma}, \texttt{BHHM}, and related
works are presently historical/method citations with no theorem consumed by
the formal proof graph.  Chua's Theorems 12.9 and 12.11 are explicitly rejected
by the target paper and are counterexample audit targets, not accepted external
inputs.  Any future use of one of these citations as a mathematical premise
must add a new row before it is accepted as a stage input or proof dependency.

\section{Explicit inputs for the selected Section 7 model}

% SOURCE: Ravenel, Complex Cobordism and Stable Homotopy Groups of Spheres,
% second edition, Theorem 3.4.5(a); conservative positive-stem specialization.
\begin{theorem}[Positive-stem sphere Adams vanishing line]
  \label{thm:external-sphere-vanishing}
  \lean{KIP126.Classical.Adams.SphereVanishingLine,KIP126.Challenge2.LiteratureInterface.sphereVanishing}
  \notready
  For the actual mod--2 Adams tower of the fixed completed sphere,
  its second page is zero in bidegree $(s,t)$ whenever
  \[
    0<t-s<2s-3.
  \]
  The zero stem is excluded, retaining the $h_0$ tower. The project uses
  this conservative specialization of Ravenel's Theorem 3.4.5(a).
  Identifying the source's spectral sequence with the fixed internal tower
  remains a model-comparison obligation. Finite CSV data cannot prove this
  all-degree vanishing statement; separation of the actual Adams filtration
  is a further internal application with convergence hypotheses.
\end{theorem}

% SOURCE: docs/STAGE0_INTERFACES.md and docs/external-inputs.json.
% The declarations below are INPUT TYPES, not proofs of the source results.
\begin{definition}[$A(M)$ on one selected model]
  \label{def:route-literature-inputs}
  \lean{KIP126.Literature.Route.Bindings,KIP126.Literature.Route.Statements,KIP126.Literature.Route.Application,KIP126.Literature.Route.Statements.toInputs,KIP126.Literature.Route.Inputs}
  \leanok
  For the fixed model $D$, the same synthetic $\eta$, and two classical tmf
  labels, the external $A(M)$ is the source record \texttt{Statements} on
  one \texttt{Bindings}. Internal source applications are delivered separately
  in \texttt{Application}. Their assembly yields the consumer \texttt{Inputs}.
  The obsolete \texttt{A = Nonempty Inputs} alias has been removed:
  \texttt{Inputs} includes internal applications and is not itself the
  external-source boundary or a field of Challenge~2.
  It has no computation table, generalized Leibniz or Mahowald conclusion,
  $C_3/C_4/C_5$, Proposition 7.8/7.9, or final-survival field.
  Defining the record does not supply an inhabitant. No default instance or
  global axiom asserts its contents.
\end{definition}

\begin{definition}[Classical and original BX inputs]
  \label{def:route-classical-bx-inputs}
  \lean{KIP126.Literature.Route.ClassicalSourceResults,KIP126.Literature.Route.ClassicalSourceBinding,KIP126.Literature.Route.ClassicalInputs,KIP126.Literature.Route.BXDistinguishedInput}
  \leanok
  Classical inputs provide an order-two class detected by the standard
  $h_5^2$, exponent two in the classical $62$-stem, its filtration gap, and
  detection of the actual maps $2$, $\eta$ and $\nu$.
  BX Proposition 7.19 and its proof are stated at one common distinguished
  $\theta_5$: the original finite criterion, the actual total-boundary
  formula, and the untruncated criterion. The LWX normalization and
  arbitrary-choice upgrade remain paper deductions.
\end{definition}

\begin{definition}[Synthetic literature and model comparisons]
  \label{def:route-synthetic-inputs}
  \lean{KIP126.Literature.Route.SyntheticInputs,KIP126.Literature.Route.RealizationInput,KIP126.Literature.Route.AlgebraInput,KIP126.Literature.Route.Applicability}
  \leanok
  The Pstragowski/BHS inputs bind lifting, Bockstein maps, rigidity,
  $E_\infty$ formulas, realization and filtration to $D$'s actual objects.
  They retain the first-quotient labels and compatibility with $\lambda$
  and restriction maps. Algebra witnesses use the existing tensor products,
  actual quotient inclusions and detector unit. Applicability separately
  requires the selected normalized $C\nu$ triangle to be distinguished and
  the classical tower condition used by Moss. These are explicit source
  application obligations, not consequences for arbitrary lift choices.
\end{definition}

\begin{definition}[May, Moss, low Toda and tmf inputs]
  \label{def:route-tools-tmf-inputs}
  \lean{KIP126.Literature.Route.MayInput,KIP126.Literature.Route.MossInput,KIP126.Literature.Route.TodaInputs,KIP126.Literature.Route.TmfInputs}
  \leanok
  May supplies signed pushpull source data; its elementwise boundary law is
  derived internally. Moss supplies a conditional
  existence statement for a detected Toda member, with defining systems,
  null products and no-crossing hypotheses retained. The low Toda inputs
  include $\lambda[h_0]=2$, $\eta^2\in\langle[h_0],\eta,[h_0]\rangle$ and
  $\lambda^2\eta\theta\in\langle2,\theta,2\rangle$ when $2\theta=0$.
  The low ring facts and the two secondary-operation calculations are separate:
  the latter require actual synthetic Toda comparison evidence, not just the
  C-motivic formula in IWX. They do not discard the high-stem indeterminacy.
  The tmf fields are
  classical Hurewicz results and a small low-filtration vanishing region;
  the synthetic local injectivity used in Proposition 7.8 is still derived
  from these inputs and the computation tables.
\end{definition}

\section{Published differentials, extensions, and stable-operation results}

\begin{proposition}[Generic Toda cosets and juggling]
  \label{prop:generic-toda-coset-juggling}
  \lean{KIP126.Foundation.TodaInterface,KIP126.StableHomotopy.Toda.indeterminacy_complete,KIP126.StableHomotopy.Toda.relation_iff_indeterminacy,KIP126.StableHomotopy.Toda.juggling,KIP126.Def.Solution.todaInterface}
  \uses{def:toda-bracket}
  \leanok
  For the cone-defined Toda relation in a triangulated category, a bracket
  exists exactly when the two successive composites vanish. Given a member
  $x_0\in\langle f,g,h\rangle$, its entire indeterminacy is
  \[
    x\in\langle f,g,h\rangle\quad\Longleftrightarrow\quad
    x-x_0=(\Sigma f)y+zh
  \]
  for suitable $y:\Sigma Y\to W$ and $z:\Sigma X\to Z$.
  If $hd=0$, there is $y\in\langle g,h,d\rangle$ with
  $xd=-(\Sigma f)y$. The additional triangulated hypothesis and sign are
  explicit; this does not calculate the named brackets below.
\end{proposition}
\begin{proof}
  The migrated proofs use exactness of distinguished triangles, completion
  of triangle morphisms and the octahedral axiom. The reverse coset inclusion
  follows from the two already proved indeterminacy actions. No historical
  global axiom is imported.
  \leanok
\end{proof}

% SOURCE: the cone-based Toda relation and its exact-functor transport.
\begin{theorem}[Toda naturality, suspension, products and shuffle]
  \label{thm:toda-product-identities}
  \uses{def:toda-bracket,prop:generic-toda-coset-juggling}
  \lean{KIP126.Foundation.TodaNaturalityInterface,KIP126.Foundation.TodaFunctorInterface,KIP126.Foundation.TodaTensorInterface,KIP126.Def.Solution.todaNaturalityInterface,KIP126.Def.Solution.todaFunctorInterface,KIP126.Def.Solution.todaTensorInterface}
  \notready
  The same cone-defined Toda relation satisfies pre- and postcomposition
  containments, both intermediate-factor absorption containments, suspension
  with the sign $(-1)^n$ and the specified shift comparison, and transport
  under an exact functor. In particular, exact left and right tensor functors
  give product containments. If $fg=gh=hd=0$, then for each $z$,
  \[
    z\in\langle f,g,h\rangle d
      \quad\Longleftrightarrow\quad
    z\in-(\Sigma f)\langle g,h,d\rangle.
  \]
  The image sets retain their complete indeterminacy. These nine precise
  statements and three Solution delivery theorems are now present; their
  proofs remain placeholders. Specific Section~7 bracket values and proofs
  that their indeterminacy vanishes remain Main obligations.
\end{theorem}

% SOURCE: MainPaper/main.tex:140--150, specialized to j=4.
% VERIFIED-IN: KIP126/Classical/Adams/Basic.lean, declaration
% KIP126.Classical.Adams.adamsOneLineDifferentials_h₄_degrees_bound.
\begin{theorem}[$h_4$ $d_2$ tracer bullet]
  \label{thm:external-adams-h4-d2-slice}
  \uses{def:classical-adams-e2-e3-shape-slice}
  \lean{KIP126.Classical.Adams.adamsOneLineDifferentials_h₄_degrees_bound}
  \leanok
  For a chosen spectrum-bound sphere $E_2$ presentation with its lawful
  bilinear, unital, associative, Leibniz-compatible and nondegenerate
  internal multiplication, the located \texttt{adams\_one\_line} result is
  consumed together with the convergence-and-detection witnesses and supplies the
  relation
  \[
    d_2(h_4)=h_0h_3^2.
  \]
  The source is in bidegree $(1,16)$ and the target in $(3,17)$.  The consumer
  retains the explicit source locator carried by the result; it does not turn
  the literature result into a project axiom.
\end{theorem}

\begin{proof}
  The exact consumer extracts the located result, retains the chosen
  spectrum's convergence-and-detection data, checks all lawful algebra fields, and
  derives the displayed bidegrees from the actual Mathlib $E_2$ differential
  component and the supplied sphere product.
\end{proof}

% SOURCE: MainPaper/main.tex:140-150; citations Adams/Mahowald--Tangora/May.
% VERIFIED-IN: MainPaper/main.tex:140-150 contains a typo: survival is j <= 3, while d_2 is nonzero for j >= 4.
\begin{theorem}[Adams one-line differentials]
  \label{thm:external-adams-one-line}
  \uses{def:adams-hj}
  \lean{KIP126.Challenge2.AdamsOneLineInterface,KIP126.Interface.Solution.adamsOneLine,KIP126.Interface.Solution.adamsOneLine_d2,KIP126.Interface.Solution.adamsHi_nonzeroSurvival_iff}
  \notready
  The internal statement fixes the one-line as a copy of $\mathbb F_2$
  generated by the specified $h_j$ in degrees $(1,2^j)$, and zero in all
  other internal degrees. It states that $h_0,h_1,h_2,h_3$ survive nontrivially and
  that, for $j\ge4$,
  \[
    d_2(h_j)=h_0h_{j-1}^2\ne0.
  \]
  This is the mathematically consistent reading of the two consecutive
  sentences in \texttt{MainPaper/main.tex:140--150}; the printed ``$j\ge3$
  survives'' is recorded as a source typo, not adopted as a project theorem.
  The same interface states May's nonzero permanence of $h_0h_2$, $h_0h_3$,
  $h_2h_4$ and $h_j^2$ for $0\le j\le3$, located at MainPaper:157--159.
  The classes are actual cobar cup products transported by the fixed internal
  comparison; all new Solution proofs remain placeholders.
\end{theorem}

\begin{proof}
  This node is a literature input.  Instantiate the external-result record at
  the cited theorem and transport its grading to
  Definition~\ref{def:adams-bidegree}; no internal proof of the classical
  calculation is claimed.
\end{proof}

% SOURCE: MainPaper/main.tex:2791, citing Bruner--Rognes (BR21).
% COORDINATES: Zenodo record 14875701, v126.3.cw49, kervaire_csv.rar;
% tmf_AdamsE2_generators.csv and tmf_AdamsE2_relations.csv (13 generators, 72 relations).
\begin{theorem}[tmf BR21 differential and coordinate compatibility (am14)]
  \label{thm:external-tmf-br21-d3-slice}
  \uses{def:hf2-adams-tower,def:adams-bidegree,def:unital-multiplication-data}
  \lean{KIP126.Classical.Adams.Tmf.CsvE2.v2Sixteen,KIP126.Classical.Adams.Tmf.CsvE2.betaFiveG,KIP126.Challenge2.TmfDifferentialInterface,KIP126.Interface.Solution.tmfDifferentialInterface}
  \lean{KIP126.Classical.Adams.Tmf.CsvE2.one,KIP126.Classical.Adams.Tmf.CsvE2.mul,KIP126.Classical.Adams.Tmf.firstMultiplication,KIP126.Classical.Adams.Tmf.unitSecondCycle}
  \lean{KIP126.Classical.Adams.Tmf.E2Presentation.RespectsMultiplication,KIP126.Classical.Adams.Tmf.E2Presentation.RespectsUnit,KIP126.Challenge2.StandardTmfMultiplicativeInterface,KIP126.Interface.Solution.tmfMultiplicativeInterface}
  \notready
  This delivery slice asks for an algebra object $T$ and an additive
  coordinate comparison, for $t\le261$, from the fixed polynomial quotient
  over $\mathbb F_2$ with all 13 CSV generators and all 72 relations to the
  actual $H\mathbb F_2$-Adams $E_2$ of $T$. In that quotient, the notation
  $v_2^{16}$ denotes $w_2^2$, the square of generator 12 in degree $(8,56)$;
  $\beta^5g$ uses generators 7 and 9 in degrees $(3,18)$ and $(4,24)$.
  Both labels are transported by the same comparison into the same tower,
  where the requested Bruner--Rognes equation is
  \[
    d_3(v_2^{16})=\beta^5g,
    \qquad (16,112)\longrightarrow(19,114).
  \]
  The algebra object's actual unit $S\to T$ fixes its classical, Adams-page
  and synthetic Hurewicz maps; these are not independently chosen maps.
  The strengthened delivery requires the same comparison $\phi$ to preserve
  that unit and the actual first-layer multiplication. The latter is formed
  from $T\wedge T\to T$, multiplication of the same $H\mathbb F_2$
  coefficients, and the fixed smash-stage and cofiber-layer comparisons,
  under the chosen tensor exactness and shift compatibilities. For every
  pair of second-cycle representatives $a,b$ with
  $[a]=\phi_{s,t}(x)$ and $[b]=\phi_{s',t'}(y)$, it requires
  \[
    [\mu_1(a,b)]=\phi_{s+s',t+t'}(xy)
       \quad\text{when }t+t'\le261.
  \]
  Here $xy$ is multiplication in the fixed CSV quotient, restricted to its
  homogeneous parts; no independent page product is supplied. The image
  of its unit is the second-page class obtained by projecting the actual
  unit $S\to T$ from Adams stage zero.
  The differential field uses common-cycle representatives on $E_3$ and asserts
  only this differential equality, without nonzero survival or detection.
  Multiplicative compatibility is now a precise condition, not a proved
  comparison. The preservation of homogeneous parts and of second cycles
  also retains explicit proof placeholders in the lower property modules.
  The geometric identification of $T$ with tmf, the vanishing of the image of $\theta_5$,
  and the required 125-stem detection remain separate obligations.
  MainPaper:2791 attributes the equation to power operations in BR21;
  neither that attribution nor the fixed CSV quotient proves it internally.
  The Solution delivery proofs remain placeholders.
\end{theorem}

% SOURCE: MainPaper/main.tex:253-263; bibliography key Adams58.  The local
% source inventory does not contain a checked copy of the original convergence
% proof, so this remains a provenance-carrying input rather than a ready node.
\begin{theorem}[Classical Adams convergence and detection input]
  \label{thm:external-classical-adams-convergence}
  \uses{def:admissible-spectrum,def:ss-convergence-detection-witness}
  \notready
  For a 2-complete connective spectrum of finite type, an explicit external
  result supplies a convergence-and-detection witness for the mod--$2$ Adams
  spectral sequence: its algebraic $E_\infty$ page is the associated graded of
  the 2-complete stable homotopy filtration, and its nonzero permanent classes
  detect the corresponding filtered homotopy elements.  The project records
  this as the standing hypothesis used in
  \texttt{MainPaper/main.tex:253--263}, with the classical source cited there
  (bibliography key \texttt{Adams58}); it does not prove convergence for
  arbitrary unbounded filtered complexes.
\end{theorem}

\begin{proof}
  Instantiate the source-bearing external-result record at the admissible
  spectrum and transport its page and filtration conventions to the concrete
  classical Adams sequence.  The source locator and the witness remain part of
  the hypothesis consumed by later definitions.
\end{proof}

% SOURCE: Ravenel, Complex Cobordism, Theorem 2.2.14; cited at MainPaper/main.tex:273-275.
\begin{theorem}[Factorization criterion for map filtration]
  \label{thm:external-map-filtration-factorization}
  \uses{def:adams-filtration,def:hf2-homology}
  \lean{KIP126.Classical.Adams.AdamsFiltrationAtLeast.hasMod2ZeroFactorization}
  \leanok
  Fix the actual Adams tower of a specified mod-two Eilenberg--Mac Lane
  object $H$, an associative multiplication with the same unit, and
  preservation of zero morphisms by $H\wedge-$.  For every positive integer
  $k$, if $\operatorname{AF}(f)\ge k$, then $f$ is a composite of exactly
  $k$ composable maps, each inducing zero on $H$-homology in every degree.
  The finite chain records its actual intermediate objects and arrows.
  The restriction $k>0$ is explicit: an empty chain represents an identity,
  whereas every map has Adams filtration at least zero.
  This forward implication is established from the constructed tower;
  Ravenel Theorem 2.2.14 remains the published source, not an additional
  hypothesis of this proof.  No converse is asserted here.
\end{theorem}

\begin{proof}
  \leanok
  The specified multiplication splits the $H$-smash of each Adams unit.
  Consequently every successor map in the actual Adams tower induces zero
  on $H$-homology.  Compose the lift witnessing
  $\operatorname{AF}(f)\ge k$ with the first of the $k$ successor maps;
  this composite still induces zero on homology.  The remaining $k-1$
  successor maps give the required chain, whose composite is $f$ by the
  tower composition law.  No literature proof or project axiom is used
  in this implication.
\end{proof}

% SOURCE: May author PDF AddJan01.pdf, TC3 pp.12--13, Lemma 4.6 p.14.
\begin{definition}[Signed May source data]
  \label{thm:external-may-smash-boundary}
  \uses{def:smash-triangle}
  \lean{KIP126.Literature.Route.MayContext,KIP126.Literature.Route.MayPushpullData,KIP126.Literature.Route.MaySourceResults}
  \leanok
  Fix the left and right tensor suspension comparisons and exactness witnesses.
  The source input gives the TC3 pushpull vertex, its three actual maps and
  the compatible-pair lifting property of Lemma 4.6. The two boundary paths
  differ by a minus sign. This records the part of TC3 used here, not a
  definition or construction of its entire axiom system. Supplying this
  source input on the selected model remains unproved.
\end{definition}

\begin{theorem}[Conditional signed boundary and sign removal]
  \label{thm:route-may-signed-adapter}
  \uses{thm:external-may-smash-boundary}
  \lean{KIP126.Interface.Solution.Literature.Route.may_signed_boundary_of_source,KIP126.Interface.Solution.Literature.Route.may_boundary_projected_of_exponent_two}
  \leanok
  Given the stated source data and a compatible pair $(a,b)$, there is $c$
  with the required second square and $\partial_U a=-\partial_T c$.
  For an additive projection $q$ satisfying $q(x)+q(x)=0$ for every $x\in\pi_{n-1}(X\wedge X')$, the projected boundaries are equal. No exponent-two property of
  actual homotopy groups follows merely from mod-two Adams $E_2$ coordinates.
\end{theorem}
\begin{proof}
  \leanok
  Lift $(a,b)$ to the source vertex, apply its third map, and compose the
  signed source identity. Apply additivity and the displayed exponent-two
  premise for the unsigned projected statement.
\end{proof}

\begin{definition}[Toda source facts and internal secondary comparison]
  \label{def:route-toda-source-application}
  \uses{def:route-tools-tmf-inputs,def:toda-bracket}
  \lean{KIP126.Literature.Route.TodaSourceData,KIP126.Literature.Route.TodaSourceResults,KIP126.Literature.Route.TodaSecondaryComparison,KIP126.Literature.Route.TodaApplication}
  \leanok
  Source facts retain one $[h_0]$, its first-quotient label, low ring relations
  and low indeterminacy. Internal comparison separately requires the actual
  low Toda member, a symmetric Toda star operation, and its identification
  with the $\lambda^2\eta$ product. IWX Corollary 6.2 is C-motivic and does
  not itself deliver these facts on the chosen synthetic model.
\end{definition}
\begin{theorem}[Conditional Toda assembly]
  \label{thm:route-toda-adapter}
  \uses{def:route-toda-source-application}
  \lean{KIP126.Interface.Solution.Literature.Route.todaApplication_of_secondary,KIP126.Interface.Solution.Literature.Route.toda_of_source}
  \leanok
  The specified secondary comparison gives the two local Toda memberships.
  Together with the source facts it produces the existing Toda consumer
  record, preserving the same $[h_0]$ and all precise premises.
\end{theorem}
\begin{proof}
  \leanok
  Substitute the star/product identity into the supplied symmetric-bracket
  theorem and assemble the fields. This does not construct the comparison.
\end{proof}

\begin{definition}[Full localization and completion comparison]
  \label{def:route-kernel-source-binding}
  \uses{def:route-synthetic-inputs}
  \lean{KIP126.Literature.Route.RealizationKernelSourceData,KIP126.Literature.Route.RealizationKernelSourceResults,KIP126.Literature.Route.RealizationKernelBinding}
  \leanok
  The full source category carries its actual $\lambda$ localization.
  Completion is left adjoint to the fully faithful inclusion of the chosen
  complete model; source spheres are identified with model spheres only
  after completion. The source finite-power kernel theorem is restricted
  to the inclusion of the selected object closure. Its model binding
  explicitly compares realization-zero and all $\lambda$ powers.
  Compactness of a hypercomplete sphere is not assumed.
\end{definition}
\begin{theorem}[Conditional finite-power kernel transport]
  \label{thm:route-kernel-adapter}
  \uses{def:route-kernel-source-binding}
  \lean{KIP126.Interface.Solution.Literature.Route.realizationKernel_of_source}
  \leanok
  Given the source kernel theorem and binding, a class on a selected object
  realizes to zero exactly when a finite $\lambda$ power kills it.
  This does not assert one-step $\lambda$ injectivity or a particular torsion bound.
\end{theorem}
\begin{proof}
  \leanok
  Use the actual adjunction equivalence on representatives, the
  realization-zero comparison, and injectivity of that equivalence together
  with its $\lambda$ compatibility.
\end{proof}

\begin{definition}[Moss on the fixed classical source convergence]
  \label{def:route-moss-source}
  \uses{def:route-classical-bx-inputs,def:route-tools-tmf-inputs}
  \lean{KIP126.Literature.Route.MossSourceInput}
  \leanok
  Moss's source statement is formulated on the same classical sphere
  convergence used by the other classical inputs. It retains three
  detections, two null products, a Massey defining system, both no-crossing
  premises and residual injectivity. The source assertion itself is not
  proved here; the original Moss scan has not been independently verified.
\end{definition}
\begin{theorem}[Moss convergence transport]
  \label{thm:route-moss-source-adapter}
  \uses{def:route-moss-source}
  \lean{KIP126.Interface.Solution.Literature.Route.mossInputOfClassicalSource}
  \leanok
  Given the source statement and the equality identifying its convergence
  with the route's sphere convergence, the existing route Moss input follows.
\end{theorem}
\begin{proof}
  \leanok
  Rewrite along the convergence equality. No defining-system, crossing,
  residual or indeterminacy hypothesis is discarded.
\end{proof}

\section{Published synthetic and comparison results}

% SOURCE: Pstragowski, Synthetic Spectra and the Cellular Motivic Category, category construction and synthetic analogue functor; summarized at MainPaper/main.tex:755-758.
\begin{theorem}[Existence and functoriality of the synthetic context]
  \label{thm:external-synthetic-foundation}
  \uses{def:synthetic-context}
  \notready
  An explicit Pstragowski external result supplies the synthetic category and
  functor of Definition~\ref{def:synthetic-context}, including stability,
  presentability, the symmetric monoidal structure, lax monoidality of $\nu$,
  and preservation of filtered colimits.
\end{theorem}

\begin{proof}
  Store the cited construction and the synthetic-analogue functor as a
  source-bearing external result, then package only the listed operations and
  coherences into the project interface.  No claim that $\nu$ preserves all
  cofibers is included.
\end{proof}

% SOURCE: Pstragowski, Corollary 4.45
% (label cor:ctau_is_an_algebra; Source/Pst/source/synthetic_spectra.tex:2247-2253).
% VERIFIED-IN: BHSmot, Appendix B, Lemma prop:mod-fil
% (Source/BHSmot/source/Filtered.tex:95-105).
\begin{theorem}[$\lambda$-quotient ring structure]
  \label{thm:external-lambda-quotient-ring}
  \uses{def:synthetic-lambda-quotient,def:lambda-quotient-operadic-target}
  \notready
  The cofiber $S^{0,0}/\lambda$ has a unique commutative, equivalently
  $E_\infty$, algebra structure compatible with the canonical quotient map
  $S^{0,0}\to S^{0,0}/\lambda$.
\end{theorem}

\begin{proof}
  Instantiate Pstragowski's located corollary; the filtered-object construction
  in BHSmot gives an independent checked copy for $C\tau$.  This external
  theorem concerns the first cofiber only and supplies no assertion about
  $S^{0,0}/\lambda^n$ for $n>1$.
  Pstragowski identifies the specified quotient arrow with the
  $0$-truncation of the connective unit. Transferring its uniqueness to the
  concrete relative unit moduli above still requires the synthetic model
  and under-unit comparisons. A strict operad-map fiber, or a claim of
  actual topological contractibility, cannot replace those obligations.
\end{proof}

% SOURCE: Pstragowski, theorem
% thm:tau_invertible_synthetic_spectra_are_just_spectra
% (Source/Pst/source/synthetic_spectra.tex:2355-2360), together with
% Proposition prop:spectral_yoneda_embedding_the_tau_inversion_of_the_synthetic_analogue
% (Source/Pst/source/synthetic_spectra.tex:2344-2346).
% VERIFIED-IN: Source/BHS/source/SynRevBigraded.tex:34-40.
\begin{theorem}[$\lambda$-inversion recovers classical spectra]
  \label{thm:external-lambda-inversion}
  \uses{def:synthetic-context,def:synthetic-lambda-quotient}
  \notready
  The spectral Yoneda embedding identifies the category of all spectra
  $\operatorname{Sp}$ with the full subcategory of $\lambda$-invertible
  $H\mathbb F_2$-synthetic spectra, and this equivalence is canonically
  symmetric monoidal.  For every spectrum $X$, the canonical map from $\nu X$
  to its spectral Yoneda image is a $\lambda$-inversion; hence
  $\lambda^{-1}\nu X\simeq X$ naturally.  No completion restriction is part of
  this generic-fiber statement.
\end{theorem}

\begin{proof}
  Apply the two located Pstragowski results and retain their monoidal and
  naturality data.  The separate restriction to hypercomplete synthetic
  spectra and $H\mathbb F_2$-local spectra occurs only later in Pstragowski
  (\texttt{Source/Pst/source/synthetic\_spectra.tex:2867--2874}) and is not
  used here.  Transport of a later map or triangle must cite this node and its
  relevant compatibility explicitly.
\end{proof}

% SOURCE: Pstragowski Lemma 4.23; MainPaper/main.tex:759-767, Proposition prop:1f7950df.
\begin{theorem}[Cofiber criterion for $\nu$]
  \label{thm:external-nu-cofiber-criterion}
  \uses{def:synthetic-context,def:cofiber-triangle,def:hf2-homology}
  \notready
  For a cofiber sequence $X\to Y\to Z$, the sequence
  $\nu X\to\nu Y\to\nu Z$ is a synthetic cofiber sequence if and only if
  \[
    0\to H_*(X;\mathbb F_2)\to H_*(Y;\mathbb F_2)
      \to H_*(Z;\mathbb F_2)\to0
  \]
  is short exact.
\end{theorem}

\begin{proof}
  This is the specialization of Pstragowski's Lemma 4.23 carried by an
  external-result parameter.  Check that the project's homology and cofiber
  conventions match the cited statement, then transport it into the abstract
  contexts.
\end{proof}

% SOURCE: BHS Theorem A.8; MainPaper/main.tex:799-807, Theorem thm:rigid.
\begin{theorem}[Synthetic Adams rigidity]
  \label{thm:external-synthetic-rigidity}
  \uses{def:classical-adams-ss,def:synthetic-adams-ss,def:bhs-aim-regrading,def:synthetic-lambda-quotient}
  \notready
  For every admissible $X$, a BHS external result gives
  \[
    {}^{\mathrm{syn}}E_2^{*,*,*}(\nu X)
      \cong E_2^{*,*}(X)\otimes\mathbb F_2[\lambda].
  \]
  Every classical differential $d_r(x)=y$ corresponds to the synthetic
  differential $d_r(x)=\lambda^{r-1}y$, and all synthetic Adams differentials
  arise this way, after applying the regrading adapter.
\end{theorem}

\begin{proof}
  Instantiate BHS Theorem A.8, transport it along
  Definition~\ref{def:bhs-aim-regrading}, and verify the three coordinates of
  each differential.  The resulting isomorphism and differential
  correspondence remain an explicit external-result parameter.
\end{proof}

% SOURCE: BHS SynRevBigraded.tex:149--159, thm:synthetic-Adams (1a)--(1c),
% retaining dfn:E-complete and dfn:strong-conv at 123--136;
% MainPaper/main.tex:813--823, thm:17e90ac0.
\begin{theorem}[Finite $\lambda$-Bockstein lifting and differential input]
  \label{thm:external-finite-lambda-bockstein}
  \lean{KIP126.Challenge2.FiniteLambdaBocksteinInterface}
  \uses{def:first-lambda-quotient-comparison,def:finite-lambda-bockstein-maps,def:actual-adams-strong-convergence}
  \notready
  Fix the same $H\mathbb F_2$ unit, $\nu$, first-quotient comparison $Q$,
  and suspension/cofiber coherence. Assume that $X$ is $H\mathbb F_2$-nilpotent
  complete and that its actual Adams tower is strongly convergent in the
  sense of Definition~\ref{def:actual-adams-strong-convergence}.
  For every $q\ge1$ and $x\in E_2^{s,t}(X)$, the finite part of BHS
  Theorem A.1 asserts
  \[
    x\in Z_q^{s,t}(X)\quad\Longleftrightarrow\quad
    \exists z\in\pi_{t-s,t}(\nu X/\lambda^q),\qquad
      \rho_{1,q,*}z=e_{s,t}^{-1}(x).
  \]
  Here $Z_q$ is the paper cycle level: it means that $x$ represents a class
  on $E_{q+1}$. Thus $q=1$ imposes no earlier differential condition, while
  $q=2$ means $d_2x=0$.
  For such an $x$ there is a suitable lift $z$ satisfying the same restriction
  equation and the actual page relation
  \[
    d_{q+1}([x])=[-b_q(z)]
      \in E_{q+1}^{s+q+1,t+q}(X).
  \]
  The Lean statement uses \texttt{DifferentialAt}; it retains the
  $B_q$ representative indeterminacy rather than asserting a strict,
  lift-independent equality in $E_2$. The negative sign is the source's
  sign. Removing it requires the mod-two property on the compared first
  quotient, not an exponent-two assertion about all synthetic homotopy groups.
\end{theorem}
\begin{proof}
  Supply the source-bearing finite input for this same model and comparison.
  The external manifest records its provenance; the Lean input remains an
  explicit mathematical parameter. This does not prove the parameter or the
  full Bockstein spectral-sequence comparison.
\end{proof}

% SOURCE: BHS SynRevAdams.tex:300--304, the displayed beta_r/E_(r+1)
% identification; BHS Theorem A.1; MainPaper/main.tex:813--823.
% MainPaper line 819's tau denotes the deformation map called lambda here.
\begin{theorem}[$\lambda$-Bockstein comparison]
  \label{thm:external-lambda-bockstein}
  \uses{thm:external-synthetic-rigidity,thm:external-finite-lambda-bockstein,def:lambda-normalized-sequence,def:pagewise-sequence-comparison,def:bhs-aim-regrading,def:synthetic-lambda-quotient}
  \notready
  Under the preceding source hypotheses, for $q\ge1$ the intrinsic $\lambda$-Bockstein
  page $q$ corresponds to synthetic Adams page $q+1$, compatibly with
  differentials and passage to the next page. In tower indices $(a,n,w)$,
  the AIM regrading is $(s,t,w)=(w+a-n,w+a,w)$.
  In particular the intrinsic first differential corresponds to synthetic
  $d_2$. Definition~\ref{def:lambda-normalized-sequence} now constructs this
  degree and page reindexing using the actual $E_1$-based tower sequence.
  The first-quotient and suitable finite-lift parts are stated
  separately above. The ``$\tau$'' in
  \texttt{MainPaper/main.tex:819} is corrected to $\lambda$ in this context.
  The raw and regraded sequences are defined using pending property proofs.
  Their comparison with the specified synthetic Adams family, the initial
  page and shifted quotient-layer identifications, and the finite rigidity
  model comparison remain unfinished. The required comparison must commute
  with differentials and successor homology maps; agreement of displayed
  pages does not identify arbitrary raw ambient SSData objects.
  Definition~\ref{def:pagewise-sequence-comparison} now supplies this generic
  comparison language, but not its first-quotient and $\nu$ representative
  bindings in the specified model.
\end{theorem}

\begin{proof}
  Use the constructed residual-tower pages, retaining the first differential,
  and establish their pending laws. Transport the located BHS comparison
  to the same synthetic model through the displayed degree and page shifts.
  Check the initial map against the
  specified first-quotient comparison and the actual cofiber boundary.
  These model comparisons remain separate obligations from the finite input
  and from defining the regraded sequence.
\end{proof}

% SOURCE: BHS Corollary A.9; MainPaper/main.tex:841-847, Proposition prop:30e8b746.
\begin{theorem}[$E_\infty$ of $\nu X$]
  \label{thm:external-synthetic-einfty-nu}
  \uses{thm:external-synthetic-rigidity,def:ss-elementwise-page-presentation,def:ss-einfty-presentation}
  \notready
  In AIM grading,
  \[
    E_\infty^{s,t,w}(\nu X)\cong
    \begin{cases}
      Z_\infty^{s,t}(X)/B_{1+t-w}^{s,t}(X),&t\ge w,\\
      0,&t<w.
    \end{cases}
  \]
  This is the regraded form of BHS Corollary A.9.
\end{theorem}

\begin{proof}
  Instantiate the external corollary, apply the BHS-to-AIM regrading, and
  check that the Bockstein filtration is $t-w$.  For nonpositive boundary
  indices use the explicit zero convention, not an implicit natural-number
  subtraction.
\end{proof}

% SOURCE: BHS Corollary A.11; MainPaper/main.tex:849-855, Proposition prop:59f111f.
\begin{theorem}[$E_\infty$ of a finite $\lambda$-quotient]
  \label{thm:external-synthetic-einfty-quotient}
  \uses{thm:external-lambda-bockstein,def:ss-elementwise-page-presentation,def:ss-einfty-presentation}
  \notready
  For $r\ge2$,
  \[
    E_\infty^{s,t,w}(\nu X/\lambda^r)\cong
    \begin{cases}
      Z_{r-t+w}^{s,t}(X)/B_{1+t-w}^{s,t}(X),&0\le t-w<r,\\
      0,&\text{otherwise}.
    \end{cases}
  \]
  This is BHS Corollary A.11 after regrading.
\end{theorem}

\begin{proof}
  Transport the cited result through the grading adapter and identify its
  truncation length with the exponent $r$.  Prove separately that every
  nonpositive cycle or boundary subscript denotes the zero group, so no
  truncated subtraction is hidden in the statement.
\end{proof}

% SOURCE: BHS Lemma 9.15; MainPaper/main.tex:912-919, Proposition prop:ef21f9bc.
\begin{theorem}[Synthetic lift of a filtered map]
  \label{thm:external-synthetic-lift}
  \uses{def:adams-filtration,def:synthetic-context,def:synthetic-lambda-quotient}
  \notready
  If $\operatorname{AF}(f)=k<\infty$, an explicit BHS result supplies a map
  $\widetilde f:\Sigma^{0,k}\nu X\to\nu Y$ with
  $\nu f=\lambda^k\widetilde f$.  This is an arbitrary chosen full lift; the
  result does not make it a component of a lifted distinguished triangle and
  does not assign a cofiber equivalence to it.
\end{theorem}

\begin{proof}
  Instantiate BHS Lemma 9.15 at $f$, transport its suspension convention, and
  retain the chosen factorization and its equality as fields of the external
  result.
\end{proof}

% SOURCE: proof of BHS Lemma 9.15; MainPaper/main.tex:929-940, Proposition prop:41561db2.
\begin{theorem}[Lift of a classical distinguished triangle]
  \label{thm:external-synthetic-triangle-lift}
  \uses{thm:external-nu-cofiber-criterion,thm:external-synthetic-lift,def:synthetic-lambda-quotient}
  \notready
  For a distinguished triangle
  $X\xrightarrow fY\xrightarrow gZ\xrightarrow h\Sigma X$ with
  $\operatorname{AF}(h)>0$ and short exact $H\mathbb F_2$-homology, an
  explicit BHS input supplies a synthetic distinguished triangle ending in
  $\Sigma^{1,0}\nu X$ and a particular triangle-compatible map $\widehat h$
  satisfying $\nu h=\lambda\widehat h$.  If
  $\operatorname{AF}(h)=k<\infty$ and a full lift $\widetilde h$ is chosen
  separately, the source identifies $\widehat h$ with
  $\lambda^{k-1}\widetilde h$ only modulo a $\lambda$-torsion map, not by an
  equality of maps.
\end{theorem}

\begin{proof}
  Use the cofiber criterion on the first two maps and retain the located
  construction in the proof of BHS Lemma 9.15 as the chosen connecting map.
  Verify the suspension
  $\Sigma^{0,-1}\nu\Sigma X=\Sigma^{1,0}\nu X$ in AIM grading.  The final
  torsion qualification is exactly the comparison recorded in
  \texttt{MainPaper/main.tex:942--943}.
\end{proof}

% SOURCE: BHS Proposition A.13; MainPaper/main.tex:1016-1019.
\begin{theorem}[$\lambda$-adic completeness]
  \label{thm:external-synthetic-lambda-complete}
  \uses{def:lambda-rho-delta-triangle}
  \lean{KIP126.StableHomotopy.SequentialHomotopyLimit,KIP126.Classical.Adams.IsENilpotentComplete,KIP126.Synthetic.Context.IsLambdaComplete,KIP126.Synthetic.Context.LambdaAdicCompletion,KIP126.Synthetic.LambdaAdicCompletenessInterface,KIP126.Def.Solution.nu_lambdaAdicCompletion}
  \notready
  For an $H\mathbb F_2$-nilpotent-complete spectrum $X$, an explicit BHS
  external result supplies a natural equivalence
  \[
    \nu X\simeq\varprojlim_m\nu X/\lambda^m
  \]
  compatible with the quotient maps and the finite
  $\lambda$--$\rho$--$\delta$ triangles.
  The completeness condition on $X$ is essential. The new Lean condition is
  acyclicity of the actual residual Adams tower, expressed by invertibility
  of $1-\mathrm{shift}$ on its countable product. The homotopy limit is the
  fiber in the Milnor triangle, retaining derived-limit information; its
  projections are the actual quotient inclusions of one specified finite
  tower. The interface also fixes the residual boundary squares and the
  equivalence with $\lambda$-completeness. These are precise statements;
  construction of the chosen model, its tower and the comparison is unproved.
\end{theorem}

\begin{proof}
  Instantiate BHS Proposition A.13 and transport its grading and inverse-system
  maps to the conventions of Definition~\ref{def:lambda-rho-delta-triangle}.
  The archived \texttt{Source/BHS/source/SynRevAdams.tex:137--148} identifies
  $E$-nilpotent completeness of $X$ with $\tau$-completeness of $\nu X$;
  it does not assert completeness for arbitrary $X$.
\end{proof}

\section{Published Kervaire and geometric results}

% SOURCE NOTE: Barratt--Jones--Mahowald inductive method and Burklund--Xu
% Proposition 7.19; MainPaper/main.tex:2117--2125, whose quantifier is ``for
% some theta_5''.  The Burklund--Xu proof internally uses the commutative
% algebra tower of Construction 7.7, but that proof dependency is represented
% only by this external-result source note and does not enter the project graph.
\begin{theorem}[Original BJM/BX criterion for the source choice]
  \label{thm:external-bjm-bx-criterion}
  \lean{KIP126.Kervaire.BJMOriginalCriterion}
  \lean{KIP126.Kervaire.BJMUntruncatedCriterion}
  \uses{def:synthetic-theta5-choice}
  \notready
  Burklund--Xu Proposition 7.19 states the finite criterion using
  $\eta(\theta_5^{\mathrm{src}})^2=0$ modulo $\lambda^r$ for survival of
  $h_6^2$ to $E_{r+3}$. In the paper's grading this lies in
  $\pi_{125,130}(S^{0,0}/\lambda^r)$.
  The proof also identifies $\delta_1(h_6^2)=\lambda\eta(\theta_5^{\mathrm{src}})^2$
  and gives the untruncated permanence criterion.
  The finite predicate now uses the actual internal tower and standard Milnor
  class, the first-quotient detection of the specified eta and theta choices,
  and the actual quotient map. Its mathematical input requires an explicit proof.
  Canonical comparison, geometric identification and the cited theorem's
  model-specific witness remain pending. The separate untruncated predicate
  uses the same standard nonzero-survival target; it is not a proved theorem.
\end{theorem}

\begin{proof}
  Source: Burklund--Xu, \texttt{kervairev2.tex}, lines 587--600.
  The finite $\lambda\eta/\lambda^{r+1}$ formulation is a separate paper
  deduction, Proposition~\ref{prop:bjm-bx-normalized-criterion}; it must not
  be supplied by labelling that strengthened interface as this external input.
\end{proof}

% SOURCE: Burklund--Xu Theorem 7.15, proof (the paragraph containing property
% (*), the quadratic-construction diagram, and Q_0(a), Q_0(b)); Proposition
% 7.19, proof (the attaching-map formula before naturality);
% MainPaper/main.tex:2133-2139, Remark rem:theta5choice.
\begin{theorem}[B--X quadratic-construction cell calculus]
  \label{thm:external-bx-quadratic-cell-calculus}
  \uses{def:stable-context,def:synthetic-context,def:synthetic-theta5-choice,def:lambda-rho-delta-triangle}
  \notready
  A located external-result record supplies the calculation for the
  distinguished source map used in the published proof, translated to the AIM
  grading convention.  Let $\theta_5^{\mathrm{src}}$ be the source class and
  $\widehat\theta_5^{\mathrm{src}}$ its chosen Moore extension.  Write $a$ and
  $b$ for the bottom and top cells.  In the quadratic construction:
  \begin{enumerate}
    \item the bottom quadratic cell maps by
      $Q_0(a)\mapsto(\theta_5^{\mathrm{src}})^2$;
    \item after reduction modulo $\lambda$, the top cell maps by
      $b\mapsto h_6$ and hence $Q_0(b)\mapsto h_6^2$; no comparison with a
      second Moore extension is included;
    \item after applying $\operatorname{Sq}(\widehat\theta_5^{\mathrm{src}})$,
      the bottom two cells in the resulting quadratic-construction diagram
      have the attaching-map identity
      \[
        \delta_1(Q_0(b))=\lambda\eta Q_0(a).
      \]
  \end{enumerate}
  The record contains only the source quadratic-construction map and its
  attaching calculation.  It does not contain a comparison map for a second
  choice or the pushed-forward identity for arbitrary $\theta_5$; both are
  Section~7 project obligations.
\end{theorem}

\begin{proof}
  The first two fields are the source-map statements in the proof of
  Burklund--Xu Theorem~7.15.  The third is the source-cell attaching-map
  calculation displayed in the proof of Proposition~7.19, before naturality is
  applied.  Store these source maps and their exact locators; do not attribute
  a comparison of independently chosen Moore extensions to B--X.
\end{proof}

% SOURCE: Source/BurklundXu/paper.txt:3475-3492, Proposition 7.19 in the
% archived version; MainPaper/main.tex:2133-2139, Remark rem:theta5choice.
% The universal TotalDifferentialIdentity declaration is project transport and
% is therefore not the external-result interface below.
\begin{theorem}[Total differential identity for the source choice]
  \label{thm:external-total-differential-identity}
  \lean{KIP126.Kervaire.BJMSourceTotalBoundaryIdentity}
  \lean{KIP126.Kervaire.SourceTotalDifferentialIdentity}
  \uses{def:synthetic-theta5-choice,def:lambda-rho-delta-triangle}
  \notready
  For the distinguished order-two source choice
  $\theta_5^{\mathrm{src}}$ used in the located Burklund--Xu quadratic
  construction, the input supplies the identity
  \[
    \delta_1(h_6^2)=\lambda\eta(\theta_5^{\mathrm{src}})^2
  \]
  in the $\delta_1:S^{0,0}/\lambda\to S^{1,-1}$ total-extension problem, with
  its source representative and grading conventions.  It does not quantify
  over other choices of $\theta_5$.
\end{theorem}

\begin{proof}
  Record the source bottom-cell class, the attaching-map calculation, and the
  resulting naturality comparison from Burklund--Xu Proposition 7.19, with the
  grading translation used by AIM.  The external input supplies no comparison
  between independently chosen Moore extensions.  Extending the identity to
  every order-two choice is a project obligation in
  Proposition~\ref{prop:theta5-total-differential-any-choice}.
\end{proof}

% SOURCE: Browder, The Kervaire invariant of framed manifolds and its generalization; MainPaper/main.tex:117-123, Theorem thm:browder.
\begin{theorem}[Browder criterion]
  \label{thm:external-browder-criterion}
  \lean{KIP126.Challenge2.BrowderInterface,KIP126.Challenge2.GeometryLiteratureInterface.browderCriterion}
  \uses{def:framed-kervaire-context,def:classical-adams-ss,def:adams-hj,def:ss-permanent-einfty-class}
  \notready
  Fix a stable context $\mathcal C$, a framed Kervaire context over
  $\mathcal C$, and a chosen classical Adams spectral sequence
  $A_{\mathrm{cl}}$ of $S^0_{\mathcal C}$.  An explicit Browder external result
  supplies that a framed manifold with Kervaire invariant one exists in
  dimension $n$ if and only if $n=2^{j+1}-2$ for some $j\ge1$ and the class
  $h_j^2$ has nonzero permanent survival in this same $A_{\mathrm{cl}}$.
  The root interface now uses the actual Milnor representatives on the fixed
  internal sphere tower. The geometric model and the proof of the literature
  result remain explicit inputs. The record must
  carry the precise theorem or page locator in Browder's paper and the
  Pontryagin--Thom convention; it cannot quantify over or silently select a
  second sphere Adams spectral sequence.
\end{theorem}

\begin{proof}
  Instantiate the located Browder theorem and transport its Adams grading and
  framed-cobordism notation into the two project interfaces.  This node remains
  an explicit literature parameter; no project proof of Browder's criterion is
  claimed.
\end{proof}

% SOURCE: MainPaper/main.tex:151-162 and cited classical Kervaire literature.
\begin{theorem}[Prior low-dimensional Kervaire existence]
  \label{thm:external-low-kervaire-existence}
  \lean{KIP126.Foundation.GeometryInterface,KIP126.Challenge2.GeometryLiteratureInterface,KIP126.Challenge2.GeometryLiteratureInterface.geometry}
  \uses{def:framed-kervaire-context}
  \notready
  Explicit literature results supply framed Kervaire-invariant-one manifolds
  in dimensions $2,6,14,30,$ and $62$, with one source locator for each
  dimension or a single theorem covering the complete list.
\end{theorem}

\begin{proof}
  Assemble the five located external results into a finite family indexed by
  the stated dimensions, verifying that each conclusion uses the same framed
  and Kervaire conventions as the project context.
\end{proof}

% SOURCE: Hill--Hopkins--Ravenel; cited at MainPaper/main.tex:86-92 and 151-160.
\begin{theorem}[HHR nonexistence beyond dimension $126$]
  \label{thm:external-hhr-nonexistence}
  \lean{KIP126.Challenge2.GeometryLiteratureInterface.high_nonexistence}
  \uses{def:framed-kervaire-context}
  \notready
  An explicit Hill--Hopkins--Ravenel result supplies nonexistence of
  Kervaire-invariant-one framed manifolds in dimensions
  $2^{j+1}-2$ for $j\ge7$, in the precise form and range of the cited theorem.
\end{theorem}

\begin{proof}
  Import the located HHR statement as an external-result parameter and align
  its dimension indexing and geometric conventions with the project context.
\end{proof}

% SOURCE: Barratt--Jones--Mahowald induction; MainPaper/main.tex:162.
\begin{proposition}[Low-dimensional permanence from explicit geometric inputs]
  \label{prop:internal-low-dimensional-permanence}
  \lean{KIP126.Challenge2.LowDimensionalSquarePermanence,KIP126.Interface.Solution.lowDimensionalSquarePermanence}
  \uses{thm:external-browder-criterion,thm:external-low-kervaire-existence,def:milnor-internal-hi-family}
  \notready
  Given the same geometric existence and Browder inputs, the actual internal
  classes $h_4^2$ and $h_5^2$ have nonzero permanent survival. This implication
  has a Lean proof. Its sphere classes use the fixed Def implementation,
  whose construction remains unfinished, and the geometric inputs themselves
  have not been constructed.
\end{proposition}
\begin{proof}
  Apply Browder to the supplied manifold of dimension $2^{j+1}-2$ for
  $j=4,5$. Injectivity of powers of two identifies Browder's returned index
  with $j$, giving survival of precisely the same named class.
\end{proof}

% SOURCE: Barratt--Jones--Mahowald induction; MainPaper/main.tex:162.
\begin{theorem}[BJM induction on order-two Kervaire classes]
  \label{thm:external-bjm-induction}
  \uses{def:adams-detection}
  \notready
  An explicit Barratt--Jones--Mahowald result supplies: if a class $\theta_j$
  detected by $h_j^2$ satisfies $2\theta_j=0$ and $\theta_j^2=0$, then there is
  an order-two class $\theta_{j+1}$ detected by $h_{j+1}^2$.  This result
  motivates the two open propositions but is not used to prove
  $h_6^2$ permanent.
\end{theorem}

\begin{proof}
  Store the exact inductive theorem and its range as a source-bearing external
  result.  The Blueprint records no converse and does not turn either open
  hypothesis into an assumption.
\end{proof}
